\begin{textsample}{POS Dim 7 – global\_south\_2023\_11 – Score 8.00 – global\_south\_2023\_11/103415885.txt}  \label{ex:f7_pos_300}
Arab Americans recoiled in the face of Donald Trump ' s campaign rhetoric against Muslims.
I ' ve been organising Arab Americans for almost a half cen? tury.
For the past three decades, my brother John and I have \textbf{polled} the community ' s views on a range of political topics.
I value polling because it allows you to hear and under? stand what people are saying and what their views may portend for the future.
We knew Arab Americans would be impacted by the devastation and loss of life from the ongoing dead? ly violence between Israelis and Palestinians.
To un? derstand the extent of the impact, we commissioned a \textbf{poll} to look more closely at the community ' s reac? tion to the conflict, its impact on their lives, their feel? ings about the Biden administration ' s handling of the conflict, and what their reactions might mean for the 2024 \textbf{elections}.
The results, released this week, were more striking than we could have imagined.
We found a dramatic decline in Arab American support for President Joseph Biden @ @ @ @ @ @ @ @ @ @ ued devastation of Gaza is the reason for this shift in attitudes.
When asked how they would \textbf{vote} in the 2024 \textbf{elections}, only 17 \% say they would cast a \textbf{ballot} for Biden, in marked con? trast to the 59 \% who \textbf{voted} for him in 2020.
The president ' s approval rating among Arab Americans also plummeted from 74 \% in 2020 to 29 \% in this year ' s \textbf{poll}.
The reason for this precipitous drop in support is clear : Two-thirds of Arab Americans say they have a negative view of the president ' s handling of the current violence in Palestine and Israel and two-thirds believe the U.S. should call for a cease? fire to end the hostilities.
The overall impact of the negative views toward the pres? ident and his policies shows up in a substantial drop in ex? pected \textbf{voter} support in 2024 and has a dramatic impact on party preference.
For the first time in our 26 years of \textbf{poll}? ing Arab American \textbf{voters}, a majority do not claim to prefer the \textbf{Democratic} Party @ @ @ @ @ @ @ @ @ @ outnumbered \textbf{Republicans}, two to one.
In this year ' s \textbf{poll}, 32 \% of Arab Americans identify as \textbf{Republican} com? pared to just 23 \% who identify as \textbf{Democrats}.
Steady growth in the percentage identifying as Independents is at the ex? pense of the \textbf{Democratic} Party.
The \textbf{poll} also demonstrates that Arab Americans worry about the domestic fallout from the war including the heated rheto? ric.
Eight in 10 Arab Americans are concerned that the violence will provoke anti-Arab bigotry, while two-thirds are concerned with the prospect of antisemitism.
They have high levels of concern with publicly expressing support for Palestinian rights and fear for their personal safety or acts of discrimination.
Six in 10 Arab Americans report experiencing discrimination, an increase of 6 \% since April 2023, with concern most prominent among Arab American Muslims ( 70 \% ) and Arab Americans ages 18-34 ( 74 \% ).
One-half of Arab Americans feel concerned about facing discrimination at school, work, and in their local community due to the recent violence.
Finally, the high levels @ @ @ @ @ @ @ @ @ @ president ' s policies are shared by almost all of the demographic groups covered in the \textbf{poll} -- by age, gen? der, education level, religion, and immigrant vs. native born.
In our nearly three decades of polling, only in two other mo? ments have policy issues resulted in such dramatic shifts in Arab American views -- and never in such a short period of time.
First, during the Bush years, over four years Arab Amer? icans moved decisively against the president ' s policies in Iraq and his repressive domestic agenda of violating civil liberties and his party ' s negative stereotyping of Arab Americans and American Muslims.
Second, in 2016, Arab Americans recoiled in the face of Donald Trump ' s campaign rhetoric against Mus? lims.
While Arab American attitudes toward President Biden ' s job performance have declined on par with all American vot? ers, the precipitous drop over a few weeks in support for his re? \textbf{election} and his party has been unprecedented.
Arab Americans may not be as @ @ @ @ @ @ @ @ @ @ hundreds of thousands of \textbf{voters} in Michigan, Ohio, and Pennsylvania were actively courted by the 2020 Biden campaign.
To win them back in 2024 will be an uphill climb.

% matched lemmas: ballot, democrat, democratic, election, poll, republican, vote, voter
\end{textsample}
